\section{Funções Com Parâmetros e Com Retorno}

\begin{frame}[fragile]{4. Com Parâmetros e Com Retorno}
    Este é o tipo mais completo e comum de função. Ela \textbf{recebe} dados de entrada para trabalhar, executa o processamento matemático ou lógico e \textbf{retorna} o resultado final para o programa.
    
    \vspace{0.3cm}
    \textbf{Analogia de Redes:}
    \begin{itemize}
        \item Pense em uma \textbf{Calculadora de Sub-Rede}: Você insere o endereço IP e a máscara CIDR (\textit{Classless Inter-Domain Routing}) (parâmetros), a função realiza os cálculos binários (processamento) e te entrega a quantidade de hosts válidos na rede (retorno).
        \item É como um roteador decidindo rotas: entra o IP de destino $\rightarrow$ processa na tabela $\rightarrow$ retorna a interface de saída.
    \end{itemize}
    
    \vspace{0.3cm}
    \textbf{Sintaxe:}
    \begin{lstlisting}[language=Python]
def calcular_soma(a, b):
    resultado = a + b
    return resultado
    \end{lstlisting}
\end{frame}

\begin{frame}[fragile]{Exemplo 1: Conversor de Mbps para Kbps}
    Para configurar a largura de banda (\textit{bandwidth}) de interfaces em roteadores Cisco, muitas vezes precisamos informar o valor em Kilobits por segundo (Kbps). Vamos criar uma função que converte Megabits por segundo (Mbps) para Kbps.
    
    \begin{lstlisting}[language=Python]
def converter_mbps_para_kbps(mbps):
    # Converte largura de banda para configuracao de interfaces
    kbps = mbps * 1024
    return kbps

# Configurando um link de internet de 50 Mbps:
banda_config = converter_mbps_para_kbps(50)

print(f"Comando Cisco: bandwidth {banda_config}")
# Saida: Comando Cisco: bandwidth 51200
    \end{lstlisting}
\end{frame}

\begin{frame}[fragile]{Exemplo 2: Classificador de IPs (Público vs Privado)}
    Função que recebe um IP como string e verifica se ele pertence às faixas privadas comuns de redes locais (\texttt{192.168.x.x} ou \texttt{10.x.x.x}), retornando \texttt{True} se for privado e \texttt{False} se for público.
    
    \begin{lstlisting}[language=Python]
def verificar_ip_privado(ip):
    # Verifica se o IP fornecido eh da rede privada local
    if ip.startswith("192.168.") or ip.startswith("10."):
        return True
    else:
        return False

# Testando pacotes vindos da Internet e da Intranet:
ip_origem1 = "192.168.1.50"
ip_origem2 = "8.8.8.8"

print(f"O IP {ip_origem1} eh privado? {verificar_ip_privado(ip_origem1)}")
print(f"O IP {ip_origem2} eh privado? {verificar_ip_privado(ip_origem2)}")
    \end{lstlisting}
\end{frame}
